\section{Discussion}
\label{sec:discussion}

In this section we reflect on the design decisions and trade-offs that
shaped the project.
We also compare the delivered artifact to the original proposal and
explain why they differ.

\subsection{Design Decisions}

\subsubsection{Focus on Definitions, Not an Optimizer Pass}

A major shift was to focus on proving relationships between concrete
definitions instead of implementing an optimizer function first. This made the
project more semantic and less algorithmic than originally planned. In
retrospect, that was the right choice. The hard part was not writing a toy
rewrite function. The hard part was deciding what its correctness should mean
in the presence of widths, \poison, and \vundef, and then making that notion
executable enough to support proofs.

By proving refinement directly between definitions, we could separate the
semantic core from the eventual question of how to enumerate or apply
optimization patterns. A future optimizer can be built on top of this
infrastructure instead of in parallel with it.

\subsubsection{Executable Simplification of \texttt{undef}}

Modeling \vundef was the most delicate design choice. A faithful model of full
\llvm nondeterminism is not easy.
The explicit supply model is a simplification, but it gave the project something
valuable: a concrete and inspectable operational account of how \vundef
choices are consumed. It also made it possible to formulate and prove the
width-generic \vundef theorem that is now one of the strongest examples in the
repository.

The downside is clear. The current model is not the full \llvm story. The
refinement relation is correspondingly tuned to this supply-based semantics.
This is a sound choice for the artifact as implemented, but it should not be
oversold as a definitive formalization of \llvm \vundef behavior.

\subsubsection{Major Differences from the Proposal}

The proposal described a broader subset of \llvm, including memory
operations, a small optimizer function, and proof obligations closer to
``\texttt{optimize p} is refined by \texttt{p}.'' The delivered system is narrower.
It has no memory model, no \texttt{load} or \texttt{store}, no
\texttt{select}, no multi-block control flow, and no standalone optimizer
pass.

This was not an accident or a loss of interest. It was a response to what the
project actually demanded once the implementation started. Even a tiny
\llvm-like semantics is full of choices: value representations, width
handling, scoping, the role of \poison, the shape of \vundef, and the form of
the refinement theorem. Building these pieces carefully took more effort than
the proposal made visible. We chose to finish a coherent core artifact rather
than spread the time thinly across a larger feature list.


\subsection{What Was Achieved and What Was Not}

The main achievements are:
\begin{itemize}
  \item an abstract syntax for a small \llvm-like \ir with symbolic widths;
  \item \lean quasiquoters and unexpansion support for readable examples;
  \item explicit well-formedness checks for definitions and arguments;
  \item executable denotational semantics for the implemented fragment;
  \item an explicit-supply model of \vundef;
  \item a semantic refinement relation over definitions; and
  \item machine-checked refinement theorems, including width-generic examples.
\end{itemize}

The main missing pieces relative to the proposal are equally important to say
plainly:
\begin{itemize}
  \item no memory model;
  \item no \texttt{load}, \texttt{store}, or pointer reasoning;
  \item no optimizer function that traverses programs and rewrites them;
  \item no proof that such an optimizer is globally sound;
  \item no control flow beyond a single basic block; and
  \item no fully faithful treatment of all \llvm undefined-behavior corners.
\end{itemize}

We view this as a successful partial completion rather than a failed broad one.
The delivered artifact is smaller than proposed, but it is real, executable,
and logically integrated.
